\documentclass{article}
\usepackage{amsmath}
\usepackage{hyperref}
\begin{document}
\section{Upper-Limb and Toe Capacity Source Candidates}
These primary measurements are candidate evidence, not subject-specific
limits. No registry value, verification label or feasibility gate is changed.
Coordinate signs refer only to the repository reference-pose sign map; gimbal
axes and generalized-force mappings are pose dependent.

\subsection{Forearm and Wrist}
Rezzoug, Hernandez and Gorce, \emph{Biomechanics} 1:102--117 (2021),
\href{https://doi.org/10.3390/biomechanics1010008}{doi:10.3390/biomechanics1010008},
Tables 2--3, tested seven right-handed male sports science students (age
$20.0\pm2.7$ years) with an isometric Biodex setup; the right upper limb was
instrumented but the tested side was not explicitly identified. At elbow $71.3\pm6.7^\circ$, forearm
pronation $75.1\pm7.9^\circ$, and wrist deviation $4.5\pm5.2^\circ$, mean
$\pm$ SD pronation/supination torques were $9.3\pm2.3$ / $11.8\pm5.3$ N m;
ulnar/radial torques were $17.1\pm8.0$ / $14.4\pm7.1$ N m. The corresponding
model signs are tentatively forearm negative/positive and wrist X
negative/positive. Side, posture and torque mapping remain unresolved.

Peleg et al., \emph{J. Funct. Morphol. Kinesiol.} 10:377 (2025),
\href{https://doi.org/10.3390/jfmk10040377}{doi:10.3390/jfmk10040377},
Table 1, tested 20 men and 20 women, all right-hand dominant, on the right
radiocarpal joint at $90^\circ$/s. The radial--ulnar range was 5 degrees
radial to 40 degrees ulnar, seated with the forearm strapped in neutral
pronation/supination, trunk supported and shoulders strapped. Male concentric
ulnar/radial peaks were $26.5\pm6.6$ / $14.9\pm3.4$ N m and eccentric peaks
were $30.3\pm7.1$ / $18.9\pm3.9$ N m. These data do not establish an isometric
limit or complete velocity-dependent envelope.

\subsection{Shoulder, Scapula and Toes}
Nicolozakes et al., \emph{Med. Sci. Sports Exerc.} 53:2354--2362 (2021),
\href{https://doi.org/10.1249/MSS.0000000000002708}{doi:10.1249/MSS.0000000000002708},
Table 1, tested the dominant right arm of eight men and nine women at 90 degrees
shoulder abduction, 20 degrees horizontal adduction and 90 degrees elbow
flexion, with scapula clamped. Male horizontal adduction/abduction was
$68\pm24$ / $45\pm16$ N m at neutral rotation and $42\pm10$ / $38\pm13$ N m
at 90 degrees external rotation. Reference-pose X shoulder signs are
plausibly related, but posture, stabilization and gimbal mapping remain open.
Rezzoug reports shoulder flexion/extension $71.6\pm17.0$ / $87.5\pm25.7$ N m
and internal/external rotation $53.9\pm12.0$ / $41.7\pm5.8$ N m; Y and Z model
coordinate-force mappings need review and these are not horizontal-axis data.
Williams et al., \emph{N. Am. J. Sports Phys. Ther.} 4:200--209 (2009),
\href{https://pmc.ncbi.nlm.nih.gov/articles/PMC2953350/}{full text}, measured
protraction/retraction forces in 30 healthy women, not independent scapular
joint torques. Hile et al., \emph{Sensors} 23:4654 (2023),
\href{https://doi.org/10.3390/s23104654}{doi:10.3390/s23104654}, Table 4,
measured extension/flexion forces of the hallux in 38 adults (18 men), seated
with 90 degrees knee flexion and neutral ankle. Hallux force is not the model's
all-toe MTP torque; no moment arm or lesser-toe contribution was established.

These citations do not authorize changing any assumption-marked capacity.
Full source, measurement side, joint frames, signs, posture, population and
dynamic transfer must be reviewed before any proposed registry update.
\end{document}
